\documentclass[11pt]{article}

\usepackage[margin=1in]{geometry}
\usepackage{amsmath}
\usepackage{amssymb}
\usepackage{amsthm}
\usepackage{booktabs}
\usepackage{parskip}

\newtheorem{conjecture}{Conjecture}
\newtheorem{proposition}{Proposition}
\theoremstyle{remark}
\newtheorem*{remark}{Remark}

\newcommand{\GL}{\mathrm{GL}}
\newcommand{\Fq}{\mathbb{F}_q}
\newcommand{\LC}{\mathrm{LC}}
\DeclareMathOperator{\rad}{rad}

\title{Disproof of Conjecture 00000001192\\[2pt]
\large The stated leading coefficient is not an integer}
\author{}
\date{}

\begin{document}
\maketitle

\section{The conjecture}

\begin{conjecture}[as filed]\label{conj}
Let $c(n,q)$ denote the number of subgroups of $\GL_n(\Fq)$ whose order equals
that of the commutator (derived) subgroup. Then $c(n,q)$ is an
\emph{integer-coefficient} polynomial in $q$ of degree $n^2-n$ with leading
coefficient
\begin{equation}\label{eq:LC}
  \LC(n) \;=\; \frac{\prod_{p \mid n} p}{n!\, n^{n}} .
\end{equation}
\end{conjecture}

\medskip
\noindent\textbf{Verdict: FALSE.} The two structural assertions in
Conjecture~\ref{conj} --- that the polynomial has \emph{integer coefficients}
and that its leading coefficient is $\LC(n)$ as in \eqref{eq:LC} --- are
mutually inconsistent. We make no claim about the actual value of $c(n,q)$;
the refutation concerns only the stated kind of polynomial and the stated
leading coefficient.

\section{The claimed leading coefficient}

Reducing \eqref{eq:LC} to lowest terms, write $\LC(n) = a_n / b_n$. Since
$\rad(n) := \prod_{p\mid n} p$ divides $n$, and $n \mid n^n \mid n!\,n^n$, we
always have $\rad(n) \mid n!\,n^n$; hence $a_n = 1$ and
$b_n = n!\,n^n / \rad(n)$ for every $n \ge 1$ (with $\rad(1)=1$). The first
values are recorded in Table~\ref{tab:lc}.

\begin{table}[h]
\centering
\begin{tabular}{@{}rrrr@{}}
\toprule
$n$ & degree $d=n^2-n$ & $\LC(n)$ & integer?\\
\midrule
$1$ & $0$  & $1$              & yes\\
$2$ & $2$  & $1/4$            & no\\
$3$ & $6$  & $1/54$           & no\\
$4$ & $12$ & $1/3072$         & no\\
$5$ & $20$ & $1/75000$        & no\\
$6$ & $30$ & $1/5598720$      & no\\
\bottomrule
\end{tabular}
\caption{The claimed leading coefficient $\LC(n)$, exactly. For every $n\ge2$
the reduced numerator is $1$ and the reduced denominator exceeds $1$, so
$\LC(n)$ is a non-integer rational. The degree $d=n^2-n\ge2$ for all $n\ge2$.}
\label{tab:lc}
\end{table}

For instance,
\[
\LC(2)=\frac{2}{2!\cdot 2^{2}}=\frac{2}{8}=\frac14,\qquad
\LC(3)=\frac{3}{3!\cdot 3^{3}}=\frac{3}{162}=\frac{1}{54},\qquad
\LC(4)=\frac{2}{4!\cdot 4^{4}}=\frac{2}{6144}=\frac{1}{3072}.
\]

\begin{proposition}\label{prop:nonint}
For every integer $n\ge 2$, $\LC(n)$ is not an integer.
\end{proposition}

\begin{proof}
As noted, $\rad(n)\mid n$, so $\rad(n)\mid n^n$, and therefore
$\rad(n) \mid n!\,n^n$. Thus the reduced numerator of
$\LC(n)=\rad(n)/(n!\,n^n)$ is $1$, while the reduced denominator is
$n!\,n^n/\rad(n)$. For $n\ge 2$ we have $\rad(n)\le n$, hence
\[
\frac{n!\,n^n}{\rad(n)} \;\ge\; \frac{n!\,n^n}{n} \;=\; n!\,n^{\,n-1}
\;\ge\; 2\cdot 1 \;=\; 2 \;>\; 1,
\]
using $n!\ge 2$ and $n^{n-1}\ge 1$. A rational with reduced numerator $1$ and
reduced denominator $>1$ is not an integer.
\end{proof}

\section{Why this refutes the conjecture}

\subsection{The integer-coefficient reading}

Suppose, as Conjecture~\ref{conj} asserts, that $c(n,q)$ is a polynomial
\[
c(n,q) = a_d q^d + a_{d-1}q^{d-1} + \dots + a_0,\qquad a_j\in\mathbb{Z},
\]
of degree $d=n^2-n$. The leading coefficient $a_d$ is then an integer, because
it is one of the integer coefficients of the polynomial. But the conjecture also
identifies $a_d$ with $\LC(n)$. Taking $n=2$ (so $d=2$ and $\LC(2)=1/4$), this
would require the integer $a_2$ to equal $1/4$, i.e.\ there would have to exist
$a\in\mathbb{Z}$ with
\[
4a = 1 .
\]
No such integer exists. The same contradiction occurs at $n=3$ ($a_6=1/54$),
$n=4$ ($a_{12}=1/3072$), and in fact at every $n\ge2$ by
Proposition~\ref{prop:nonint}. Consequently no polynomial with integer
coefficients of the stated degree has the stated leading coefficient:
Conjecture~\ref{conj} is false as written.

\subsection{The charitable ``integer-valued'' reading}

One might weaken ``integer-coefficient'' to ``integer-valued'' (i.e.\
$c(n,q)\in\mathbb{Z}$ for all integers $q$), since a PORC formula often only
needs to take integer values. This does not rescue the conjecture. A classical
fact about integer-valued polynomials is:

\begin{proposition}\label{prop:fact}
If $P\in\mathbb{Q}[x]$ is integer-valued and has degree $d$, then
$d!\,\cdot\,(\text{leading coefficient of }P)\in\mathbb{Z}$.
\end{proposition}

\begin{proof}[Sketch]
The binomial polynomials $\binom{x}{k}$ ($k=0,\dots,d$) form a $\mathbb{Z}$-basis
of the integer-valued polynomials of degree $\le d$. In this basis
$P=\sum_{k=0}^d b_k\binom{x}{k}$ with $b_k\in\mathbb{Z}$, and
$\binom{x}{k}$ has leading coefficient $1/k!$. Hence the leading coefficient of
$P$ is $b_d/d!$, so $d!$ times it is $b_d\in\mathbb{Z}$.
\end{proof}

Applying Proposition~\ref{prop:fact} to the conjectured degree $d=n^2-n$:
\begin{itemize}
  \item $n=2$: $d=2$ and
    $\displaystyle 2!\,\LC(2)=2\cdot\frac14=\frac12\notin\mathbb{Z}$;
  \item $n=3$: $d=6$ and
    $\displaystyle 6!\,\LC(3)=720\cdot\frac{1}{54}=\frac{720}{54}=\frac{40}{3}\notin\mathbb{Z}$.
\end{itemize}
So even under the weaker integer-valued reading, the claimed leading coefficient
is impossible at the two smallest cases $n=2$ and $n=3$. The ambiguity between
the two readings does not rescue the conjecture.

\subsection{Arithmetic content, in divisibility form}

The two obstructions used above are exactly the following elementary facts,
which are the ones formalised in the accompanying Lean development:
\[
\nexists\, m\in\mathbb{N}\ \bigl(m\cdot 4 = 1\bigr),\qquad
\nexists\, m\in\mathbb{N}\ \bigl(m\cdot 2 = 1\bigr),\qquad
\nexists\, a\in\mathbb{Z}\ \bigl(4a=1\bigr),
\]
and, for the integer-valued reading at $n=3$,
\[
\nexists\, m\in\mathbb{N}\ \bigl(m\cdot 54=720\bigr)
\qquad\Bigl(\tfrac{720}{54}=\tfrac{40}{3}\Bigr).
\]

\section{Scope, and the PORC caveat}

The quantity $c(n,q)$ itself --- the number of subgroups of $\GL_n(\Fq)$ whose
order equals that of the derived subgroup --- is not otherwise analysed in this
submission. The refutation is entirely about the \emph{stated leading
coefficient} being arithmetically impossible for the \emph{stated kind of
polynomial} (integer-coefficient, degree $n^2-n$). This is a self-contained
inconsistency in the text of the conjecture, independent of any group theory.

If the author intended a PORC description, in which the coefficient is a
constituent depending on $q=p^f$ or on the residue class of $q$, or intended a
different normalisation (for example, a degree-$0$ or lower-degree polynomial,
or a coefficient such as $\rad(n)/(n!\,n^n)$ viewed only after multiplication by
a further factor), then the formula as written is still self-inconsistent: as
written it demands an integer-coefficient polynomial whose leading coefficient
is the non-integer rational $\LC(n)$. To make the statement consistent one must
change at least one of: the class of polynomials (allowing rational
coefficients), the degree, or the normalisation of the leading coefficient. No
such qualification appears in the conjecture as filed.

\begin{remark}
Because the contradiction already occurs at $n=2$ (degree $2$), the refutation
does not depend on any subtlety of Higman's PORC phenomenon or on the actual
count $c(n,q)$; it is a statement about rational arithmetic.
\end{remark}

\section*{Reproducing}

The exact values of $\LC(n)$ for $n=1,\dots,12$ (and the $d!\,\LC(n)$
integer-valued test) are produced by
\begin{center}
\texttt{python3 reproduce.py}
\end{center}
which exits $0$ on \textsc{pass}. The finite arithmetic facts above are
formalised in core Lean~4 under \texttt{lean4/} (no Mathlib, no \texttt{sorry});
build with \texttt{lake build} and audit with \texttt{lake env lean Check.lean}.

\end{document}
